% Apunts de teoria en PDF · 1r ESO · Unitat 1: Nombres naturals
% Mateix estil que els apunts de Batxillerat (franja i caixes de teoria en blau, exemples en roig), però més visuals,
% curts i amb dibuixos, per a l'alumnat de 3r ESO. Es compila amb pdflatex dues vegades (per als enllaços de l'índex),
% amb logo-fentmates.pdf (_eines/logo/pdf/horitzontal-clar.pdf) al costat.
\documentclass[11pt,a4paper]{article}
\usepackage[utf8]{inputenc}
\usepackage[T1]{fontenc}
\usepackage[catalan]{babel}
\usepackage[sfdefault,book,lining]{FiraSans}
\usepackage[italic,noplusnominus,noequal]{mathastext}
\linespread{1.18}
\usepackage[a4paper,top=2.4cm,bottom=2.2cm,left=1.6cm,right=1.6cm,headheight=1.1cm,headsep=0.5cm,footskip=1cm]{geometry}
\usepackage{amsmath,amssymb,cancel}
\usepackage{icomma}
\usepackage{graphicx}
\usepackage[table]{xcolor}
\usepackage{tikz}\usetikzlibrary{calc,positioning,arrows.meta,decorations.pathreplacing,shapes.geometric}
\usepackage{fontawesome5}
\usepackage{ccicons}
\usepackage{fancyhdr}
\usepackage{enumitem}
\usepackage{array}
\usepackage{needspace}
\usepackage[most]{tcolorbox}
\usepackage{titlesec}
\usepackage[hidelinks]{hyperref}
\hypersetup{pdftitle={Apunts de teoria. Unitat 1: Nombres naturals. 1r ESO},pdfsubject={Matemàtiques 1r ESO},pdfkeywords={nombres naturals, sistema decimal, aproximació, truncament, arrodoniment, operacions, propietats, jerarquia, problemes, potències, potències de base 10, propietats de les potències, arrel quadrada}}

% colors de la web (bloc d'Àlgebra: blau)
\definecolor{fmnavy}{HTML}{1B2748}
\definecolor{fmblue}{HTML}{3D6AA0}
\definecolor{fmbluel}{HTML}{4E81BD}
\definecolor{fmyellow}{HTML}{FFC94D}
\definecolor{fmgold}{HTML}{E9A919}
\definecolor{fmaccent}{HTML}{FF7A50}
\definecolor{fmdef}{HTML}{1D4F73}\definecolor{fmdefbg}{HTML}{EEF4F9}
\definecolor{fmex}{HTML}{A3203A}\definecolor{fmexbg}{HTML}{FBF0F2}
\definecolor{fmsf}{HTML}{C2410C}
\definecolor{fmok}{HTML}{15803D}
\definecolor{fmmuted}{HTML}{5D6685}
\definecolor{fmsurface}{HTML}{EEF1F8}
\definecolor{fmA}{HTML}{2F6FD6}\definecolor{fmB}{HTML}{EA6A10}\definecolor{fmR}{HTML}{E11D48}\definecolor{fmK}{HTML}{16A34A}\definecolor{fmV}{HTML}{7C3AED}

% lletres
\newcommand{\comic}{\fontfamily{ComicNeue-TLF}\selectfont}
\newcommand{\heavy}{\fontseries{eb}\selectfont}
\newcommand{\semi}{\fontseries{sb}\selectfont}
\DeclareUnicodeCharacter{2212}{-}

\setlength{\parindent}{0pt}
\setlength{\parskip}{7pt}
\setlength{\lineskip}{4pt}\setlength{\lineskiplimit}{4pt}
\setlength{\emergencystretch}{3em}
\setlist{itemsep=3pt,topsep=3pt,leftmargin=1.4em}
\setcounter{secnumdepth}{1}
\renewcommand{\thesection}{1.\arabic{section}}
\titleformat{\section}[hang]{\color{fmnavy}\fontsize{23}{27}\heavy\selectfont}{\tcbox[on line,colback=fmblue,colframe=fmblue,arc=3mm,boxsep=2pt,left=6pt,right=6pt,top=3pt,bottom=3pt]{\color{white}\thesection}}{0.45em}{}[\vspace{5pt}{\color{fmblue}\titlerule[2pt]}]
\titlespacing*{\section}{0pt}{4pt}{6pt}
\titleformat{\subsection}{\color{fmblue}\Large\semi}{}{0pt}{\faCaretRight\ }
\titlespacing*{\subsection}{0pt}{8pt}{2pt}

% frase de cada apartat, «a mà»
\newcommand{\frase}[1]{\par\vspace{-2pt}{\comic\Large\color{fmbluel}\faCommentDots[regular]\ \ #1\par}\vspace{2pt}}
\newcommand{\seccio}[3]{\Needspace{16\baselineskip}\section{#2}\hypertarget{#1}{}\frase{#3}}
% notes a mà dins dels exemples
\newcommand{\nota}[1]{{\comic\small\color{fmex!85!black}#1}}
\newcommand{\notab}[1]{{\comic\small\color{fmblue}#1}}
% bé / malament
\newcommand{\be}[1]{{\color{fmok}\faCheckCircle}\ \ensuremath{#1}}
\newcommand{\mal}[1]{{\color{fmR}\faTimesCircle}\ {\color{fmR}\ensuremath{#1}}}
% pas numerat
\newcommand{\pas}[1]{\tikz[baseline=(p.base)]\node[circle,fill=fmblue,text=white,inner sep=1.2pt,minimum size=1.25em,font=\small\heavy](p){#1};}
% etiqueta de color
\newcommand{\chip}[2]{\tcbox[on line,colback=#1!14,colframe=#1!60,arc=2mm,boxsep=1pt,left=5pt,right=5pt,top=1pt,bottom=1pt]{#2}}
\newcommand{\fmkey}[1]{\colorbox{fmyellow!55}{\strut #1}}
\newcommand{\cA}[1]{{\color{fmA}#1}}\newcommand{\cB}[1]{{\color{fmB}#1}}\newcommand{\cK}[1]{{\color{fmK}#1}}\newcommand{\cR}[1]{{\color{fmR}#1}}\newcommand{\cV}[1]{{\color{fmV}#1}}

% caixes
\tcbset{fmbase/.style={enhanced,arc=3mm,boxrule=1.2pt,left=4mm,right=4mm,top=2mm,bottom=2.5mm,before skip=8pt,after skip=8pt,fonttitle=\heavy\large,before upper=\raggedright}}
\newtcolorbox{idea}[1][Idea clau]{fmbase,colback=fmdefbg,colframe=fmdef,coltitle=white,title={\faKey\ \ #1}}
\newtcolorbox{exemple}[1][Exemple]{fmbase,colback=fmexbg,colframe=fmex,colbacktitle=fmexbg,coltitle=fmex,title={\faPencil*\ \ #1},titlerule=0.5pt,titlerule style={fmex!40,dashed}}
\newtcolorbox{compte}[1][Compte!]{fmbase,colback=fmsf!6,colframe=fmsf,coltitle=white,fonttitle=\comic\bfseries\Large,title={\faExclamationTriangle\ \ #1}}
\newtcolorbox{truc}[1][Truc]{fmbase,colback=fmyellow!20,colframe=fmgold,colbacktitle=fmyellow,coltitle=fmnavy,fonttitle=\comic\bfseries\Large,title={\faLightbulb[regular]\ \ #1}}

% dibuix de la targeta de la unitat, en blanc
\newcommand{\fmcard}{\begin{tikzpicture}[x=0.062cm,y=-0.062cm,white,line width=3pt,line cap=round]
  \begin{scope}[transparency group,opacity=0.15]
  \node[anchor=base west,inner sep=0pt,font=\fontsize{30}{30}\selectfont] at (14,44) {$2^5\quad\sqrt{49}$};
  \draw[line width=2.4pt] (20,95) -- (230,95);
  \foreach \x in {30,60,90,120,150,180,210} \draw[line width=2.4pt] (\x,87) -- (\x,103);
  \end{scope}
\end{tikzpicture}}
\newcommand{\fmbrand}{\raisebox{-0.25\height}{\includegraphics[height=0.5cm]{logo-fentmates.pdf}}}

% regla de Ruffini: #1 = a, #2 = coeficients, #3 = fila del mig (sense el primer), #4 = fila de baix, #5 = nombre de coeficients
\newcommand{\ruffini}[5]{\begin{tikzpicture}[x=1.25cm,y=0.85cm,every node/.style={font=\large}]
  \foreach \v [count=\i] in {#2} \node (t\i) at (\i,0) {$\v$};
  \foreach \v [count=\i] in {#3} {\pgfmathtruncatemacro\j{\i+1}\node (m\j) at (\j,-1) {$\v$};}
  \foreach \v [count=\i] in {#4} \node (b\i) at (\i,-2.1) {$\v$};
  \node[text=fmB] at (0,-1) {$#1$};
  \draw[line width=1.1pt,fmnavy] (0.45,0.5)--(0.45,-1.5) (0.45,-1.5)--(#5+0.5,-1.5);
  \node[draw=fmR,line width=1.3pt,rounded corners=2mm,minimum size=0.75cm] at (#5,-2.1) {};
  \draw[-{Latex[length=2.2mm]},fmK,line width=1.1pt] (t1.south) -- (b1.north);
  \pgfmathtruncatemacro\last{#5-1}
  \foreach \i in {1,...,\last} {\pgfmathtruncatemacro\j{\i+1}
    \draw[-{Latex[length=2.2mm]},fmB,line width=1.1pt,shorten >=1pt] (b\i.north east) -- (m\j.south west);
    \draw[-{Latex[length=2.2mm]},fmA!70,line width=1pt,dashed,shorten >=2pt] ([xshift=8pt]t\j.south) -- ([xshift=8pt]b\j.north);}
\end{tikzpicture}}

% capçalera i peu
\pagestyle{fancy}\fancyhf{}
\fancyhead[L]{\small\fmbrand}
\fancyhead[R]{\small\color{fmmuted}1r ESO · Unitat 1: Nombres naturals}
\fancyfoot[L]{\scriptsize\color{fmmuted}\href{https://creativecommons.org/licenses/by-nc-sa/4.0/deed.ca}{\ccbyncsa}\ \ FentMates · fentmates.org · CC BY-NC-SA 4.0}
\fancyfoot[R]{\small\color{fmmuted}\thepage}
\renewcommand{\headrulewidth}{0.4pt}\renewcommand{\headrule}{\color{fmmuted!40}\hrule height \headrulewidth}
\fancypagestyle{primera}{\fancyhf{}\renewcommand{\headrulewidth}{0pt}
  \fancyfoot[L]{\scriptsize\color{fmmuted}\href{https://creativecommons.org/licenses/by-nc-sa/4.0/deed.ca}{\ccbyncsa}\ \ FentMates · fentmates.org · CC BY-NC-SA 4.0}
  \fancyfoot[R]{\small\color{fmmuted}\thepage}}

% targeta de l'índex
\newcommand{\itile}[4]{\begin{tcolorbox}[enhanced,colback=white,colframe=fmblue!35,arc=3mm,boxrule=1pt,boxsep=0pt,left=6pt,right=6pt,top=5pt,bottom=5pt,before skip=0pt,after skip=7pt]
  \hyperlink{#1}{\begin{minipage}[c]{0.9cm}\centering{\color{fmblue}\LARGE #3}\end{minipage}\hspace{6pt}%
  \begin{minipage}[c]{0.78\linewidth}{\color{fmblue}\heavy #2}\\[-2pt]{\color{fmnavy}\semi\large #4}\end{minipage}}\end{tcolorbox}}

% taula amb vores arrodonides: el contingut es retalla amb la forma (sense @{} a les vores ni entre columnes)
\newsavebox\tabclipbox
\newcommand{\taulaclip}[3][3mm]{\sbox\tabclipbox{#3}\begin{tikzpicture}[baseline=(tcn.base)]
  \begin{scope}\clip[rounded corners=#1] (0,-\dp\tabclipbox) rectangle (\wd\tabclipbox,\ht\tabclipbox);
  \node[inner sep=0pt,outer sep=0pt,anchor=base west] (tcn) at (0,0) {\usebox\tabclipbox};\end{scope}
  \draw[#2,rounded corners=#1] (0,-\dp\tabclipbox) rectangle (\wd\tabclipbox,\ht\tabclipbox);
\end{tikzpicture}}
\newcommand{\taularod}[3][fmblue]{\taulaclip{#1,line width=1.2pt}{\renewcommand{\arraystretch}{1.5}\begin{tabular}{#2}#3\end{tabular}}}
% balança
\newcommand{\balanca}[3]{\begin{tikzpicture}[scale=#3,transform shape,every node/.style={anchor=center}]
  \fill[fmnavy] (-0.45,-1.9) rectangle (0.45,-1.75);
  \fill[fmnavy!80] (-0.12,-1.75) rectangle (0.12,0);
  \fill[fmgold] (0,0.05) circle (0.13);
  \draw[fmnavy,line width=2.4pt,line cap=round] (-2,0) -- (2,0);
  \foreach \s in {-1,1} {
    \draw[fmmuted,line width=0.8pt] (\s*2,0) -- (\s*2-0.75,-1.05) (\s*2,0) -- (\s*2+0.75,-1.05);
    \fill[fmyellow!70,draw=fmgold,line width=1pt] (\s*2-0.95,-1.05) arc[start angle=180,end angle=360,x radius=0.95,y radius=0.32] -- cycle;}
  \node[font=\Large,fill=fmA!15,draw=fmA,rounded corners=1.5mm,inner sep=3pt] at (-2,-0.68) {$#1$};
  \node[font=\Large,fill=fmB!15,draw=fmB,rounded corners=1.5mm,inner sep=3pt] at (2,-0.68) {$#2$};
\end{tikzpicture}}
% eixos amb quadrícula (dins d'un tikzpicture): xmin, xmax, ymin, ymax
\newcommand{\eixos}[4]{\draw[fmmuted!22,line width=0.4pt] (#1,#3) grid (#2,#4);
  \draw[-{Latex[length=2mm]},fmmuted,line width=0.8pt] (#1-0.2,0)--(#2+0.4,0) node[right,font=\small]{$x$};
  \draw[-{Latex[length=2mm]},fmmuted,line width=0.8pt] (0,#3-0.2)--(0,#4+0.4) node[above,font=\small]{$y$};}
% taula de traduccions (diccionari) amb les vores arrodonides
\newcommand{\dic}[2]{#1 & {\large$#2$}\\}
\newcommand{\dicratio}{0.62}
\newcommand{\dichead}{\begin{tabular}{@{}>{\raggedright}p{\dimexpr\dicratio\linewidth-12pt}>{\centering\arraybackslash}p{\dimexpr\linewidth-\dicratio\linewidth-12pt}@{}}\faComment[regular]\ Ho diem així\dots & \dots\ així\end{tabular}}
\newenvironment{dicbox}{\begin{tcolorbox}[enhanced,arc=3mm,boxrule=1.2pt,colframe=fmblue,colback=white,coltitle=white,fonttitle=\heavy,
  title={\dichead},left=0pt,right=0pt,top=0pt,bottom=1.5mm,boxsep=0pt,lefttitle=6pt,righttitle=6pt,before skip=0pt,after skip=0pt]\rowcolors{1}{white}{fmsurface}%
  \begin{tabular}{@{\hspace{6pt}}>{\raggedright}p{\dimexpr\dicratio\linewidth-12pt}>{\centering\arraybackslash}p{\dimexpr\linewidth-\dicratio\linewidth-12pt}@{\hspace{6pt}}}}
  {\end{tabular}\end{tcolorbox}}
% resum
\newcommand{\rtile}[3]{\begin{tcolorbox}[enhanced,colback=white,colframe=fmblue!45,arc=3mm,boxrule=1pt,left=3mm,right=3mm,top=2.5mm,bottom=2.5mm,fontupper=\large,
  title={\tcbox[on line,colback=fmblue,colframe=fmblue,arc=2mm,boxsep=0pt,left=4pt,right=4pt,top=1pt,bottom=1pt]{\color{white}\small\heavy #1}\ \ #2},coltitle=fmnavy,colbacktitle=fmdefbg,fonttitle=\semi,toptitle=0.5mm,bottomtitle=0.5mm]#3\end{tcolorbox}}
% portada: #1 subtítol a mà
\newcommand{\portada}[1]{\thispagestyle{primera}
\begin{tikzpicture}[remember picture,overlay]
  \fill[fmblue] (current page.north west) rectangle ([yshift=-9.4cm]current page.north east);
  \begin{scope}
    \clip (current page.north west) rectangle ([yshift=-9.4cm]current page.north east);
    \node[anchor=north west,inner sep=0pt] at ([xshift=6.6cm,yshift=-1.3cm]current page.north west) {\fmcard};
  \end{scope}
  \node[anchor=north west,inner sep=0pt] at ([xshift=1.6cm,yshift=-1.1cm]current page.north west) {\tcbox[colback=white,colframe=white,arc=2.5mm,boxsep=0pt,left=7pt,right=7pt,top=5pt,bottom=5pt]{\includegraphics[height=0.75cm]{logo-fentmates.pdf}}};
\end{tikzpicture}
\vspace*{0.55cm}
\tcbox[colback=white!15!fmblue,colframe=white!15!fmblue,arc=2.5mm,boxsep=0pt,left=6pt,right=6pt,top=3pt,bottom=3pt]{\color{white}\small\bfseries Apunts de teoria · 1r ESO · Matemàtiques}

\vspace{2pt}
{\raggedright\color{white}\fontsize{30}{36}\heavy\selectfont{\color{fmyellow}Unitat 1}\\[2pt]Nombres naturals\par}
\vspace{8pt}
{\color{white}\comic\Large #1\par}
\vspace*{\dimexpr 7.4cm-\pagetotal\relax}
}
\newcommand{\llegenda}{\par\vspace{16pt}
{\color{fmnavy}\heavy\Large\faPalette\ \ Com es llegeixen aquests apunts}\par\smallskip
\begin{tcbraster}[raster columns=4,raster equal height,raster column skip=8pt,raster before skip=12pt]
\tcbox[colback=fmdefbg,colframe=fmdef,arc=3mm,boxrule=1.2pt,halign=center,width=\linewidth]{{\color{fmdef}\Large\faKey}\\[2pt]{\color{fmdef}\heavy Idea clau}\\[-1pt]{\small el que cal saber}}
\tcbox[colback=fmexbg,colframe=fmex,arc=3mm,boxrule=1.2pt,halign=center,width=\linewidth]{{\color{fmex}\Large\faPencil*}\\[2pt]{\color{fmex}\heavy Exemple}\\[-1pt]{\small pas a pas}}
\tcbox[colback=fmsf!6,colframe=fmsf,arc=3mm,boxrule=1.2pt,halign=center,width=\linewidth]{{\color{fmsf}\Large\faExclamationTriangle}\\[2pt]{\color{fmsf}\heavy Compte!}\\[-1pt]{\small errors típics}}
\tcbox[colback=fmyellow!20,colframe=fmgold,arc=3mm,boxrule=1.2pt,halign=center,width=\linewidth]{{\color{fmgold!80!black}\Large\faLightbulb[regular]}\\[2pt]{\color{fmnavy}\heavy Truc}\\[-1pt]{\small per a recordar-ho}}
\end{tcbraster}}
\newcommand{\resumtitol}{\clearpage
\hypertarget{resum}{}
{\color{fmnavy}\fontsize{23}{27}\heavy\selectfont\faClipboardList\ \ En resum}\par\vspace{5pt}{\color{fmblue}\titlerule[2pt]}
\frase{Tota la unitat en una pàgina.}}
% els 5 passos per a resoldre problemes: #2 = text del pas 2
\newcommand{\passosproblema}[2]{\begin{center}
\begin{tikzpicture}[s/.style={draw=fmblue,fill=fmdefbg,rounded corners=3mm,line width=1.2pt,minimum width=2.75cm,minimum height=2.45cm,align=center,font=\small},
  ic/.style={font=\LARGE,text=fmblue}]
  \foreach \i/\ic/\t in {1/\faBookOpen/{\textbf{Llegim}\\i entenem},2/\faQuestionCircle/{#1},3/\faEquals/{#2},4/\faCogs/{\textbf{Resolem}},5/\faCheckDouble/{\textbf{Comprovem}\\i responem}} {
    \node[s] (s\i) at (3.3*\i,0) {};
    \node[ic] at ([yshift=-0.5cm]s\i.north) {\ic};
    \node[align=center,font=\small,anchor=north] at ([yshift=-0.95cm]s\i.north) {\t};
    \node[circle,fill=fmblue,text=white,font=\small\heavy,inner sep=1.5pt,minimum size=1.3em] at (s\i.north west) {\i};}
  \foreach \i [evaluate=\i as \j using int(\i+1)] in {1,...,4} \draw[-{Latex[length=3mm]},fmmuted,line width=1.5pt] (s\i) -- (s\j);
\end{tikzpicture}
\end{center}}

\newcommand{\lo}{\mathopen{]}}\newcommand{\ro}{\mathclose{[}}
% punts ple i buit, nombres dels eixos i taula de valors
\newcommand{\pple}[2][fmblue]{\fill[#1] #2 circle (3.6pt);}
\newcommand{\nums}[4]{\foreach \k in {#1,...,#2} {\ifnum\k=0\else\node[below,font=\scriptsize,text=fmmuted] at (\k,0) {$\k$};\fi}
  \foreach \k in {#3,...,#4} {\ifnum\k=0\else\node[left,font=\scriptsize,text=fmmuted] at (0,\k) {$\k$};\fi}}
\newcommand{\tv}[3][fmblue]{\taulaclip{#1,line width=1.2pt}{\renewcommand{\arraystretch}{1.45}%
  \begin{tabular}{>{\columncolor{#1}\color{white}\heavy}c*{12}{w{c}{0.8cm}}}#2\\\rowcolor{#1!8}\cellcolor{#1}#3\end{tabular}}}
% etiqueta d'una recta o corba
\newcommand{\etq}[3]{\node[fill=white,inner sep=1.5pt,text=#1,font=\small] at #2 {#3};}


\def\intvdot#1#2#3{\ifx#2o\filldraw[fill=white,draw=#3,line width=1.3pt] (#1,0) circle (3.3pt);\fi\ifx#2c\fill[#3] (#1,0) circle (3.5pt);\fi}
\newcommand{\intv}[5]{\begin{tikzpicture}[x=0.4cm,baseline=-3pt]
  \draw[fmmuted,line width=0.8pt] (-6,0)--(7,0);
  \foreach \k in {-5,...,6} \draw[fmmuted,line width=0.6pt] (\k,-0.08)--(\k,0.08);
  \ifx#3a\def\xa{-6}\else\def\xa{#1}\fi
  \ifx#4a\def\xb{7}\else\def\xb{#2}\fi
  \ifx#4a\draw[#5,line width=2.6pt,-{Latex[length=2.6mm]}] (\xa,0)--(\xb,0);\else\ifx#3a\draw[#5,line width=2.6pt,{Latex[length=2.6mm]}-] (\xa,0)--(\xb,0);\else\draw[#5,line width=2.6pt] (\xa,0)--(\xb,0);\fi\fi
  \intvdot{#1}{#3}{#5}\intvdot{#2}{#4}{#5}
  \ifx#3a\else\node[below=2pt,font=\scriptsize] at (#1,0) {$#1$};\fi
  \ifx#4a\else\node[below=2pt,font=\scriptsize] at (#2,0) {$#2$};\fi
\end{tikzpicture}}
\newcommand{\passos}[1]{\begin{tcolorbox}[fmbase,colback=fmdefbg,colframe=fmdef,coltitle=white,title={\faListOl\ \ Com es resolen}]#1\end{tcolorbox}}
\newcommand{\passosT}[2]{\begin{tcolorbox}[fmbase,colback=fmdefbg,colframe=fmdef,coltitle=white,title={\faListOl\ \ #1}]#2\end{tcolorbox}}
\newcommand{\barra}[3]{\filldraw[fill=#3!35,draw=#3,line width=0.9pt] (#1-0.32,0) rectangle (#1+0.32,#2);}
\newcommand{\caixa}[7]{\draw[#7,line width=1.3pt] (#2,#1) -- (#3,#1) (#5,#1) -- (#6,#1) (#2,#1-0.25) -- (#2,#1+0.25) (#6,#1-0.25) -- (#6,#1+0.25);
  \filldraw[fill=#7!20,draw=#7,line width=1.3pt] (#3,#1-0.45) rectangle (#5,#1+0.45); \draw[#7,line width=2.2pt] (#4,#1-0.45) -- (#4,#1+0.45);}
\newcommand{\recta}[2]{\draw[-{Latex[length=2.5mm]},fmnavy,line width=1.1pt] (#1-0.6,0) -- (#2+0.6,0);
  \foreach \k in {#1,...,#2} {\draw[fmnavy,line width=0.9pt] (\k,-0.12) -- (\k,0.12); \node[below=3pt,font=\small] at (\k,0) {$\k$};}}


\begin{document}
\portada{Comptar, ordenar i operar: els nombres de cada dia.}

{\color{fmnavy}\heavy\Large\faMapSigns\ \ Què aprendràs?}\par\smallskip
\begin{minipage}[t]{0.485\linewidth}
\itile{sec1}{1.1}{\faSortNumericUp}{Els nombres naturals}
\itile{sec2}{1.2}{\faIcon{cut}}{Truncament i arrodoniment}
\itile{sec3}{1.3}{\faCalculator}{Operacions amb naturals}
\itile{sec4}{1.4}{\faLightbulb}{Problemes amb naturals}
\itile{sec5}{1.5}{\faSuperscript}{Potències}
\end{minipage}\hfill
\begin{minipage}[t]{0.485\linewidth}
\itile{sec6}{1.6}{\faIcon{expand-arrows-alt}}{Potències de base 10}
\itile{sec7}{1.7}{\faListOl}{Propietats de les potències}
\itile{sec8}{1.8}{\faIcon{square-root-alt}}{Arrels quadrades}
\itile{sec9}{1.9}{\faTasks}{Problemes amb potències i arrels}
\itile{resum}{\phantom{1.9}}{\faClipboardList}{En resum}
\end{minipage}
\llegenda

% ============================================================ 1.1
\clearpage
\seccio{sec1}{Els nombres naturals}{Per a comptar i per a ordenar.}

\vfill
\begin{minipage}[c]{0.52\linewidth}
\begin{idea}[Els naturals $\mathbb N$]
{\centering\Large$\mathbb N=\{0,\ 1,\ 2,\ 3,\ 4,\ \ldots\}$\par}\smallskip
Serveixen per a \textbf{comptar} (5 caixes, 120 alumnes) i per a \textbf{ordenar} (el 3r classificat). Comencen en el 0 i no s'acaben mai: després de $n$ sempre ve $n+1$.
\end{idea}
\end{minipage}\hfill
\begin{minipage}[c]{0.45\linewidth}\centering
\begin{tikzpicture}[x=0.62cm]
  \draw[-{Latex[length=2.4mm]},fmnavy,line width=1.2pt] (0,0) -- (9.6,0);
  \foreach \k in {0,...,9} {\draw[fmnavy,line width=1pt] (\k,-0.12) -- (\k,0.12); \node[below=2pt,font=\small] at (\k,-0.1) {$\k$};}
  \foreach \k in {0,...,8} \draw[fmB,line width=1pt,-{Latex[length=1.6mm]}] (\k+0.1,0.25) to[bend left=50] (\k+0.9,0.25);
  \node[above,font=\scriptsize,text=fmB] at (4.5,0.7) {$+1$ cada pas};
\end{tikzpicture}\\[2pt]
{\comic\small\color{fmmuted}una semirecta: sempre la mateixa distància}
\end{minipage}

\vfill
\begin{minipage}[c]{0.5\linewidth}\centering
\taularod{>{\heavy}l<{\hspace{4pt}}>{\hspace{4pt}}l<{\hspace{4pt}}>{\hspace{4pt}}c<{\hspace{4pt}}}{
\rowcolor{fmblue}\color{white}Operació & \color{white}\heavy Exemple & \color{white}\heavy Dona un natural?\\
Suma & $7+5=12$ & \be{}\\
\rowcolor{fmsurface}Multiplicació & $7\cdot5=35$ & \be{}\\
Resta & $5-7=\ ?$ & \mal{} \ {\small (enters)}\\
\rowcolor{fmsurface}Divisió & $7:5=\ ?$ & \mal{} \ {\small (si no és exacta)}
}
\end{minipage}\hfill
\begin{minipage}[c]{0.47\linewidth}
\begin{tcolorbox}[fmbase,colback=fmdefbg,colframe=fmdef,coltitle=white,title={Propietat distributiva}]
{\centering$a\cdot(b\pm c)=a\cdot b\pm a\cdot c$\par}\smallskip
Per a calcular de cap: \ $7\cdot99=7\cdot(100-1)=700-7=693$\\[3pt]
Per a traure \textbf{factor comú}: \ $12\cdot15+12\cdot5=12\cdot20=240$
\end{tcolorbox}
\end{minipage}

\vfill
\begin{minipage}[c]{0.4\linewidth}\centering
\begin{tikzpicture}[font=\Large]
  \node (D) at (0,0) {$247$}; \node (d) at (1.6,0) {$12$};
  \draw[fmnavy,line width=1.2pt] (0.75,0.35) -- (0.75,-0.35) -- (2.4,-0.35);
  \node[text=fmA] at (1.6,-0.75) {$20$}; \node[text=fmR] at (0,-1.1) {$7$};
  \node[font=\scriptsize,text=fmmuted] at (-0.9,0) {dividend}; \node[font=\scriptsize,text=fmmuted] at (2.7,0) {divisor};
  \node[font=\scriptsize,text=fmA] at (2.7,-0.75) {quocient}; \node[font=\scriptsize,text=fmR] at (-0.9,-1.1) {residu};
\end{tikzpicture}
\end{minipage}\hfill
\begin{minipage}[c]{0.57\linewidth}
\begin{idea}[La divisió entera]
Si no és exacta, dona un \textbf{quocient} $q$ i un \textbf{residu} $r$, sempre \textbf{menor} que el divisor:\\[4pt]
{\centering\Large$D=d\cdot q+r$ \qquad \normalsize$247=12\cdot20+7$\par}\smallskip
És la \textbf{prova de la divisió}: si no es compleix, hi ha un error.
\end{idea}
\end{minipage}

\vfill
\begin{exemple}[Exemples]
\begin{minipage}[t]{0.48\linewidth}
250 caramels en bosses de 12: \ $250=12\cdot20+10$\\ \hspace*{1em}\boxed{20\text{ bosses i en sobren }10}\\[4pt]
$25\cdot101=25\cdot100+25=\boxed{2\,525}$
\end{minipage}\hfill
\begin{minipage}[t]{0.49\linewidth}
Quin dividend té divisor 15, quocient 8 i residu 6?\\ \hspace*{1em}$D=15\cdot8+6=\boxed{126}$\\[4pt]
$36\cdot7+36\cdot3=36\cdot10=\boxed{360}$
\end{minipage}
\end{exemple}

\vfill
\begin{compte}[El residu sempre és menor que el divisor]
$100=7\cdot13+9$ \mal{} \ no és una divisió entera: $9>7$. \ La correcta: \ $100=7\cdot14+2$ \be{} \qquad I \textbf{no es pot dividir entre 0}.
\end{compte}

\vfill

\vfill
% ============================================================ 1.2
\clearpage
\seccio{sec2}{Aproximació: truncament i arrodoniment}{Un nombre pròxim, més fàcil de llegir i de recordar.}

\vfill
\begin{idea}[Aproximar]
Canviar un nombre per un altre de pròxim i més fàcil. A un concert hi ha 44\,863 persones: el diari dirà «unes 45\,000». S'aproxima a un \textbf{ordre}: desenes, centenes, milers\dots
\end{idea}

\vfill
\begin{tcbraster}[raster columns=2,raster equal height,raster column skip=10pt]
\begin{tcolorbox}[fmbase,colback=fmA!8,colframe=fmA,coltitle=white,title={\faIcon{cut}\ \ Truncar}]
Canviem per \textbf{zeros} totes les xifres de la dreta de l'ordre: «tallem».\par\smallskip
{\centering$44\,8\underline{63}\to44\,800$ \ \nota{a les centenes}\par}
\end{tcolorbox}
\begin{tcolorbox}[fmbase,colback=fmB!8,colframe=fmB,coltitle=white,title={\faIcon{circle}\ \ Arrodonir}]
Mirem la xifra \textbf{següent}: 0, 1, 2, 3 o 4, la de l'ordre es queda igual; 5, 6, 7, 8 o 9, li \textbf{sumem 1}. Després, zeros.\par\smallskip
{\centering$44\,8\fbox{6}3\to44\,900$ \qquad $44\,8\fbox{2}3\to44\,800$\par}
\end{tcolorbox}
\end{tcbraster}

\vfill
\begin{center}
\begin{tikzpicture}[x=0.11cm]
  \draw[fmnavy,line width=1.1pt] (0,0) -- (100,0);
  \foreach \k in {0,10,...,100} \draw[fmnavy] (\k,-0.12) -- (\k,0.12);
  \node[below,font=\small] at (0,-0.1) {44\,800}; \node[below,font=\small] at (50,-0.1) {44\,850}; \node[below,font=\small] at (100,-0.1) {44\,900};
  \fill[fmR] (63,0) circle (3pt) node[above=2pt,font=\small\heavy,text=fmR] {44\,863};
  \draw[-{Latex[length=2mm]},fmK,line width=1.2pt] (63,0.5) to[bend left=20] node[above,font=\scriptsize] {més a prop} (100,0.5);
\end{tikzpicture}\\
{\comic\small\color{fmmuted}44\,863 està més a prop de 44\,900 que de 44\,800}
\end{center}

\vfill
\begin{exemple}[Exemple: \ 44\,863 aproximat de les dues maneres]
\centering\small
\taularod[fmex]{>{\heavy}l*{4}{>{\hspace{8pt}}c<{\hspace{8pt}}}}{
\rowcolor{fmex}\color{white}44\,863 & \color{white}Truncament & \color{white}Error & \color{white}Arrodoniment & \color{white}Error\\
a les desenes & 44\,860 & 3 & 44\,860 & 3\\
\rowcolor{fmexbg}a les centenes & 44\,800 & 63 & 44\,900 & \cK{37}\\
als milers & 44\,000 & 863 & 45\,000 & \cK{137}
}\par\smallskip\raggedright
L'\textbf{error} és la diferència entre l'exacte i l'aproximat: l'arrodoniment sempre dona el \textbf{més pròxim}.
\end{exemple}

\vfill
\begin{exemple}[Exemples]
\begin{minipage}[t]{0.48\linewidth}
Arrodoneix 7\,459 a les centenes: \ $\boxed{7\,500}$\\[4pt]
Arrodoneix 128\,350 als milers: \ $\boxed{128\,000}$
\end{minipage}\hfill
\begin{minipage}[t]{0.49\linewidth}
Trunca 9\,987 a les desenes: \ $\boxed{9\,980}$\\[4pt]
Arrodoneix 9\,987 a les desenes: \ $\boxed{9\,990}$
\end{minipage}
\end{exemple}

\vfill
\begin{compte}[Quan s'arrossega]
2\,996 a les centenes: la xifra següent és un 9, així que sumem 1 a les 9 centenes: es fan 10 i en portem una al miler: \ $\boxed{3\,000}$.
\end{compte}

\vfill
% ============================================================ 1.3
\clearpage
\seccio{sec3}{Operacions amb nombres naturals}{Les propietats ens ajuden a calcular més ràpid.}

\vfill
\begin{tcbraster}[raster columns=3,raster equal height,raster column skip=8pt]
\begin{tcolorbox}[fmbase,colback=fmA!8,colframe=fmA,coltitle=white,title={Commutativa},before upper=\centering]
L'ordre no importa:\par\smallskip
$7+5=5+7$\par $4\cdot9=9\cdot4$
\end{tcolorbox}
\begin{tcolorbox}[fmbase,colback=fmB!8,colframe=fmB,coltitle=white,title={Associativa},before upper=\centering]
Agrupem com vulguem:\par\smallskip
$(25\cdot4)\cdot7=25\cdot(4\cdot7)$\par $=100\cdot7=700$
\end{tcolorbox}
\begin{tcolorbox}[fmbase,colback=fmK!8,colframe=fmK,coltitle=white,title={Distributiva},before upper=\centering]
$a\cdot(b+c)=a\cdot b+a\cdot c$\par\smallskip
$7\cdot99=700-7=693$
\end{tcolorbox}
\end{tcbraster}

\vfill
\begin{compte}[La resta i la divisió no són commutatives]
$9-4\neq4-9$ \qquad $12:3\neq3:12$. \ I recorda que \textbf{no es pot dividir entre 0}.
\end{compte}

\vfill
\begin{minipage}[c]{0.42\linewidth}\centering
\begin{tikzpicture}[st/.style={draw=#1,fill=#1!12,line width=1.2pt,rounded corners=2mm,minimum height=0.8cm,anchor=south west,font=\large,inner xsep=6pt}]
  \node[st=fmV] at (0,1.8) {\pas{1}\ \ Parèntesis $(\ )\ [\ ]$};
  \node[st=fmK] at (0.5,0.9) {\pas{2}\ \ $\cdot$ \ i \ $:$};
  \node[st=fmB] at (1.0,0) {\pas{3}\ \ $+$ \ i \ $-$};
\end{tikzpicture}\\[3pt]
{\comic\small\color{fmmuted}\pas{2} i \pas{3}: d'esquerra a dreta}
\end{minipage}\hfill
\begin{minipage}[c]{0.55\linewidth}
\begin{exemple}[Exemple: \ $30-4\cdot(9-3)+12:3$]
$=30-4\cdot6+12:3$ \ \nota{\pas{1}}\\[3pt]
$=30-24+4$ \ \nota{\pas{2}}\\[3pt]
$=\boxed{10}$ \ \nota{\pas{3}}
\end{exemple}
\end{minipage}

\vfill
\begin{compte}[Errors freqüents]
\be{8+2\cdot5=18}, no $10\cdot5=50$: primer el producte. \qquad \be{12:2\cdot3=18}, no $12:6=2$: d'esquerra a dreta.
\end{compte}

\vfill
\begin{exemple}[Exemples]
\begin{minipage}[t]{0.48\linewidth}
$5+3\cdot4=5+12=\boxed{17}$\\[4pt]
$(5+3)\cdot4=8\cdot4=\boxed{32}$\\[4pt]
$40:(2\cdot4)+1=40:8+1=\boxed{6}$
\end{minipage}\hfill
\begin{minipage}[t]{0.49\linewidth}
$18-[7+(10-6)]=18-11=\boxed{7}$\\[4pt]
$12\cdot15+12\cdot5=12\cdot20=\boxed{240}$ \ \nota{factor comú}\\[4pt]
$25\cdot13\cdot4=100\cdot13=\boxed{1\,300}$
\end{minipage}
\end{exemple}

\vfill
% ============================================================ 1.4
\clearpage
\seccio{sec4}{Problemes amb nombres naturals}{Quatre passos i la paraula que diu l'operació.}

\passosproblema{\textbf{Pensa}\\un pla}{\textbf{Tria}\\les operacions}

\vfill
\begin{minipage}[c]{0.55\linewidth}\centering\small
\taularod{>{\hspace{3pt}}l<{\hspace{3pt}}>{\hspace{3pt}}c<{\hspace{3pt}}}{
\rowcolor{fmblue}\color{white}\heavy Si el problema diu\dots & \color{white}\heavy Operació\\
ajuntar, afegir, en total, entre tots & $+$\\
\rowcolor{fmsurface}llevar, gastar, quant en queda, quant falta & $-$\\
cada un\dots té\dots, vegades, el doble, el triple & $\cdot$\\
\rowcolor{fmsurface}repartir a parts iguals, quantes vegades cap & $:$
}
\end{minipage}\hfill
\begin{minipage}[c]{0.42\linewidth}
\begin{compte}[I el residu?]
De vegades cal afegir un grup més (\textbf{autobusos}: ningú es queda a terra) i de vegades no es compta (\textbf{caixes plenes}).
\end{compte}
\end{minipage}

\vfill
\begin{tcbraster}[raster columns=2,raster equal height,raster column skip=10pt,raster row skip=8pt]
\begin{exemple}[Exemple: \ autobusos]
150 alumnes van d'excursió en autobusos de 55 places. Quants autobusos calen?\\[3pt]
$150=55\cdot2+40$: sobren 40 alumnes, cal un autobús més: \ \boxed{3\text{ autobusos}}
\end{exemple}
\begin{exemple}[Exemple: \ caixes plenes]
Tenim 150 ous i els posem en caixes de 12. Quantes caixes plenes?\\[3pt]
$150=12\cdot12+6$: \ \boxed{12\text{ caixes}} \ i en sobren 6.
\end{exemple}
\end{tcbraster}

\vfill
\begin{exemple}[Exemple: \ una família compra una tele de 549 € i paga 3 mensualitats de 125 € i la resta al final. Quant paga al final?]
Mensualitats: \ $3\cdot125=375$ € \qquad Resta: \ $549-375=\boxed{174\text{ €}}$ \qquad \nota{comprova: $375+174=549$}
\end{exemple}

\vfill
\begin{exemple}[Exemple: \ en una biblioteca hi ha 12 prestatgeries de 5 lleixes amb 35 llibres cadascuna. Hi arriben 300 llibres més]
Llibres: \ $12\cdot5\cdot35=2\,100$ \qquad Ara: \ $2\,100+300=\boxed{2\,400}$ \qquad Per prestatgeria: \ $2\,400:12=\boxed{200}$
\end{exemple}

\vfill
\begin{truc}[Escriu què calcules]
En cada pas, una frase curta: «Mensualitats: 375 €». Al final, respon amb una \textbf{frase} i comprova que té sentit.
\end{truc}

\vfill
% ============================================================ 1.5
\clearpage
\seccio{sec5}{Potències}{Una multiplicació de factors iguals.}

\vfill
\begin{minipage}[c]{0.5\linewidth}\centering
\begin{tikzpicture}
  \node[font=\fontsize{40}{44}\selectfont,text=fmA] (b) at (0,0) {3};
  \node[font=\fontsize{24}{28}\selectfont,text=fmB,anchor=south west] (e) at (0.25,0.25) {4};
  \node[font=\Large] at (2.9,0) {$=3\cdot3\cdot3\cdot3=81$};
  \node[font=\small,text=fmA,draw=fmA,rounded corners=2mm,fill=fmA!8] at (-0.9,-1.1) {base};
  \node[font=\small,text=fmB,draw=fmB,rounded corners=2mm,fill=fmB!8] at (1.3,1.2) {exponent};
\end{tikzpicture}
\end{minipage}\hfill
\begin{minipage}[c]{0.47\linewidth}
\begin{idea}[Potència]
La \cA{\textbf{base}} és el factor que es repeteix; l'\cB{\textbf{exponent}}, quantes vegades.\\[3pt]
{\centering$a^n=\underbrace{a\cdot a\cdots a}_{n\text{ vegades}}$\par}
\end{idea}
\end{minipage}

\vfill
\begin{compte}[L'exponent no multiplica]
\be{2^5=2\cdot2\cdot2\cdot2\cdot2=32} \qquad \mal{2^5=2\cdot5=10}
\end{compte}

\vfill
\begin{tcbraster}[raster columns=3,raster equal height,raster column skip=8pt]
\begin{tcolorbox}[fmbase,colback=fmdefbg,colframe=fmdef,coltitle=white,title={Com es llegeixen},before upper=\centering]
$a^2$: «al quadrat»\par $a^3$: «al cub»\par $7^4$: «set elevat a quatre»
\end{tcolorbox}
\begin{tcolorbox}[fmbase,colback=fmdefbg,colframe=fmdef,coltitle=white,title={Casos especials},before upper=\centering]
$a^1=a$\par\smallskip $a^0=1$ \ (si $a\neq0$)\par\smallskip $1^n=1$
\end{tcolorbox}
\begin{tcolorbox}[fmbase,colback=fmdefbg,colframe=fmdef,coltitle=white,title={Quadrats i cubs},before upper=\centering]
\begin{tikzpicture}[scale=0.28]
  \draw[fmA,fill=fmA!20] (0,0) grid (3,3) rectangle (0,0); \node[font=\scriptsize] at (1.5,-0.7) {$3^2=9$};
  \begin{scope}[xshift=4.5cm]\draw[fmB,fill=fmB!25] (0,0) rectangle (2,2); \draw[fmB,fill=fmB!40] (0,2) -- (0.8,2.8) -- (2.8,2.8) -- (2,2) -- cycle; \draw[fmB,fill=fmB!55] (2,0) -- (2.8,0.8) -- (2.8,2.8) -- (2,2) -- cycle; \node[font=\scriptsize] at (1.4,-0.7) {$2^3=8$};\end{scope}
\end{tikzpicture}
\end{tcolorbox}
\end{tcbraster}

\vfill
\begin{exemple}[Exemples]
\begin{minipage}[t]{0.48\linewidth}
$5^3=5\cdot5\cdot5=\boxed{125}$\\[4pt]
$2^6=\boxed{64}$ \qquad $10^4=\boxed{10\,000}$\\[4pt]
$9^1=9$ \qquad $9^0=1$ \qquad $1^{20}=1$
\end{minipage}\hfill
\begin{minipage}[t]{0.49\linewidth}
En forma de potència: \ $6\cdot6\cdot6\cdot6=\boxed{6^4}$\\[4pt]
$4\cdot4\cdot4\cdot7\cdot7=\boxed{4^3\cdot7^2}$\\[4pt]
Quadrats: $1,\ 4,\ 9,\ 16,\ 25,\ 36,\ 49,\ 64,\ 81,\ 100$
\end{minipage}
\end{exemple}

\vfill
\begin{exemple}[Exemple: \ cada caixa té 6 bosses, cada bossa 6 paquets i cada paquet 6 caramels. Quants caramels hi ha en una caixa?]
$6\cdot6\cdot6=6^3=\boxed{216}$ caramels
\end{exemple}

\vfill
% ============================================================ 1.6
\clearpage
\seccio{sec6}{Potències de base 10}{Un 1 seguit de zeros.}

\vfill
\begin{center}\small
\taularod{>{\heavy}l*{7}{>{\hspace{5pt}}c<{\hspace{5pt}}}}{
\rowcolor{fmblue}\color{white}Potència & \color{white}$10^1$ & \color{white}$10^2$ & \color{white}$10^3$ & \color{white}$10^4$ & \color{white}$10^5$ & \color{white}$10^6$ & \color{white}$10^9$\\
Valor & 10 & 100 & 1\,000 & 10\,000 & 100\,000 & 1\,000\,000 & 1\,000\,000\,000\\
\rowcolor{fmsurface}Es llegeix & deu & cent & mil & deu mil & cent mil & un milió & mil milions
}
\end{center}

\vfill
\begin{idea}[El truc dels zeros]
$10^n$ és un \textbf{1 seguit de $n$ zeros}. \ Multiplicar per $10^n$ és afegir $n$ zeros: \ $37\cdot10^3=37\,000$.
\end{idea}

\vfill
\begin{minipage}[c]{0.52\linewidth}
\begin{idea}[Descomposició polinòmica]
Cada xifra val segons el lloc que ocupa: \textbf{unitats}, \textbf{desenes} ($10$), \textbf{centenes} ($10^2$), \textbf{milers} ($10^3$)\dots\\[4pt]
$4\,705=4\cdot10^3+7\cdot10^2+0\cdot10+5$
\end{idea}
\end{minipage}\hfill
\begin{minipage}[c]{0.45\linewidth}\centering
\taularod{*{4}{>{\hspace{6pt}}c<{\hspace{6pt}}}}{
\rowcolor{fmblue}\color{white}\heavy UM & \color{white}\heavy C & \color{white}\heavy D & \color{white}\heavy U\\
\Large 4 & \Large 7 & \Large 0 & \Large 5\\
\rowcolor{fmsurface}\small$4\cdot10^3$ & \small$7\cdot10^2$ & \small$0\cdot10$ & \small 5
}
\end{minipage}

\vfill
\begin{tcbraster}[raster columns=2,raster equal height,raster column skip=10pt]
\begin{tcolorbox}[fmbase,colback=fmK!8,colframe=fmK,coltitle=white,title={Nombres grans}]
La distància de la Terra al Sol: uns 150\,000\,000 km $=15\cdot10^7$ km.\par\smallskip
La població del món: uns $8\cdot10^9$ persones.
\end{tcolorbox}
\begin{tcolorbox}[fmbase,colback=fmA!8,colframe=fmA,coltitle=white,title={Comptem els zeros}]
$5\cdot10^6=5\,000\,000$\par\smallskip
$32\,000=32\cdot10^3$\par\smallskip
$10^2\cdot10^3=100\cdot1\,000=10^5$
\end{tcolorbox}
\end{tcbraster}

\vfill
\begin{exemple}[Exemples]
\begin{minipage}[t]{0.48\linewidth}
$3\,604=3\cdot10^3+6\cdot10^2+4$\\[4pt]
$2\cdot10^4+5\cdot10^2+1=\boxed{20\,501}$\\[4pt]
$7\cdot10^5=\boxed{700\,000}$
\end{minipage}\hfill
\begin{minipage}[t]{0.49\linewidth}
Un milió de persones: $10^6$\\[4pt]
$450\,000=45\cdot10^4$\\[4pt]
$10^6:10^2=10^4=10\,000$
\end{minipage}
\end{exemple}

\vfill
\begin{compte}[Els zeros del mig compten]
En $4\,705$ el 0 és la xifra de les desenes: $0\cdot10$. Si no el posem, \mal{475} és un altre nombre.
\end{compte}

\vfill
% ============================================================ 1.7
\clearpage
\seccio{sec7}{Propietats de les potències}{Amb la mateixa base, juguem amb els exponents.}

\providecommand{\prop}{}\renewcommand{\prop}[3]{\begin{tcolorbox}[fmbase,colback=fmdefbg,colframe=fmdef,coltitle=white,title={#1},before upper=\centering]
{\Large$#2$}\par\smallskip{\small$#3$}
\end{tcolorbox}}
\begin{tcbraster}[raster columns=3,raster equal height=rows,raster column skip=8pt]
\prop{\pas{1}\ Producte}{a^{m}\cdot a^{n}=a^{\cK{m+n}}}{2^{5}\cdot2^{3}=2^{8}=256}
\prop{\pas{2}\ Quocient}{a^{m}:a^{n}=a^{\cK{m-n}}}{7^{9}:7^{7}=7^{2}=49}
\prop{\pas{3}\ Potència de potència}{\left(a^{m}\right)^{n}=a^{\cK{m\cdot n}}}{\left(5^{2}\right)^{3}=5^{6}=15\,625}
\end{tcbraster}

\vfill
\begin{minipage}[c]{0.5\linewidth}
\begin{idea}[Exponent 1 i exponent 0]
$a^{1}=a$: \ $7^{1}=7$\\[4pt]
$a^{0}=1$ \ (si $a\neq0$): \ $5^{3}:5^{3}=5^{0}=1$
\end{idea}
\end{minipage}\hfill
\begin{minipage}[c]{0.47\linewidth}
\begin{truc}[Per què se sumen?]
$2^{3}\cdot2^{2}=(2\cdot2\cdot2)\cdot(2\cdot2)$: \ en total, $3+2=5$ factors iguals: \ $2^{5}$.
\end{truc}
\end{minipage}

\vfill
\begin{compte}[Només amb la mateixa base, i per a productes i quocients]
\mal{2^{2}+2^{3}=2^{5}}\ \ {\small$4+8=12$, \ però $2^5=32$}\qquad\quad
$2^{3}\cdot3^{2}$: bases diferents, calculem cada una: \ $8\cdot9=72$
\end{compte}

\vfill
\begin{exemple}[Exemples de cada propietat]
\begin{minipage}[t]{0.32\linewidth}
\pas{1} $4^{2}\cdot4^{5}=4^{7}$\\[4pt]
\hspace*{1.3em}$10\cdot10^{4}=10^{5}$\\[4pt]
\hspace*{1.3em}$2^{3}\cdot2^{3}\cdot2^{3}=2^{9}$
\end{minipage}\hfill
\begin{minipage}[t]{0.32\linewidth}
\pas{2} $3^{9}:3^{4}=3^{5}$\\[4pt]
\hspace*{1.3em}$10^{6}:10^{2}=10^{4}$\\[4pt]
\hspace*{1.3em}$8^{5}:8^{4}=8^{1}=8$
\end{minipage}\hfill
\begin{minipage}[t]{0.32\linewidth}
\pas{3} $\left(3^{2}\right)^{4}=3^{8}$\\[4pt]
\hspace*{1.3em}$\left(10^{3}\right)^{2}=10^{6}$\\[4pt]
\hspace*{1.3em}$\left(6^{5}\right)^{1}=6^{5}$
\end{minipage}
\end{exemple}

\vfill
\clearpage
\subsection{Apliquem les propietats}
\begin{exemple}[Exemples: \ escriu com una sola potència]
\begin{minipage}[t]{0.48\linewidth}
$3^{4}\cdot3^{2}\cdot3=3^{4+2+1}=\boxed{3^{7}}$ \ \nota{$3=3^1$}\\[5pt]
$5^{8}:5^{3}=5^{8-3}=\boxed{5^{5}}$\\[5pt]
$\left(2^{3}\right)^{4}=2^{3\cdot4}=\boxed{2^{12}}$
\end{minipage}\hfill
\begin{minipage}[t]{0.49\linewidth}
$10^{2}\cdot10^{3}\cdot10^{4}=\boxed{10^{9}}$\\[5pt]
$6^{10}:6^{9}=6^{1}=\boxed{6}$\\[5pt]
$\left(4^{2}\right)^{0}=4^{0}=\boxed{1}$
\end{minipage}
\end{exemple}

\vfill
\begin{tcbraster}[raster columns=2,raster equal height,raster column skip=10pt]
\begin{exemple}[Exemple: \ $\left(2^{5}\cdot2^{3}\right):2^{6}$]
$=2^{8}:2^{6}$ \ \nota{\pas{1} primer el parèntesi}\\[3pt]
$=2^{2}=\boxed{4}$ \ \nota{\pas{2}}
\end{exemple}
\begin{exemple}[Exemple: \ $\left(7^{2}\right)^{3}:7^{4}$]
$=7^{6}:7^{4}$ \ \nota{\pas{3}}\\[3pt]
$=7^{2}=\boxed{49}$ \ \nota{\pas{2}}
\end{exemple}
\begin{exemple}[{Exemple: \ $\left[\left(3^{2}\right)^{3}\cdot3^{4}\right]:3^{8}$}]
$=\left[3^{6}\cdot3^{4}\right]:3^{8}=3^{10}:3^{8}$\\[3pt]
$=3^{2}=\boxed{9}$
\end{exemple}
\begin{exemple}[Exemple: \ $\left(5^{4}\cdot5^{2}\right):\left(5^{3}\right)^{2}$]
$=5^{6}:5^{6}$\\[3pt]
$=5^{0}=\boxed{1}$ \ \nota{dividend i divisor iguals}
\end{exemple}
\end{tcbraster}

\vfill
\begin{exemple}[Exemples: \ completa l'exponent que falta]
$4^{\square}\cdot4^{3}=4^{7}$: \ $\square=\boxed{4}$ \qquad $9^{8}:9^{\square}=9^{5}$: \ $\square=\boxed{3}$ \qquad $\left(2^{\square}\right)^{3}=2^{12}$: \ $\square=\boxed{4}$ \qquad $\dfrac{10^{7}\cdot10^{2}}{10^{6}}=10^{\square}$: \ $\square=\boxed{3}$
\end{exemple}

\vfill
\begin{compte}[Errors típics]
$3^{2}\cdot3^{4}=3^{6}$, \ \mal{9^{6}}: la base no es multiplica\\[3pt] $\left(2^{3}\right)^{2}=2^{6}$, \ \mal{2^{5}}: els exponents es multipliquen\\[3pt] $5^{6}:5^{2}=5^{4}$, \ \mal{5^{3}}: es resten, no es divideixen
\end{compte}

\vfill
% ============================================================ 1.8
\clearpage
\seccio{sec8}{Arrels quadrades}{L'operació inversa d'elevar al quadrat.}

\vfill
\begin{minipage}[c]{0.55\linewidth}
\begin{idea}[Arrel quadrada]
{\centering\Large$\sqrt a=b\iff b^2=a$\par}\smallskip
$\sqrt{49}=7$ perquè $7^2=49$. \ El signe $\sqrt{\ }$ és el \textbf{radical} i el nombre de dins, el \textbf{radicand}.\\[3pt]
Elevar al quadrat i fer l'arrel es \textbf{desfan} l'una a l'altra.
\end{idea}
\end{minipage}\hfill
\begin{minipage}[c]{0.42\linewidth}\centering
\begin{tikzpicture}[scale=0.32]
  \draw[fmA,fill=fmA!20] (0,0) grid (7,7); \draw[fmA,line width=1.3pt] (0,0) rectangle (7,7);
  \node[below,font=\small] at (3.5,0) {7 cm}; \node[font=\small\heavy] at (3.5,3.5) {49 cm$^2$};
\end{tikzpicture}\\
{\comic\small\color{fmmuted}un quadrat de 49 cm$^2$ fa 7 cm de costat}
\end{minipage}

\vfill
\begin{center}\small
\taularod{>{\heavy}l*{15}{>{\hspace{1pt}}c<{\hspace{1pt}}}}{
\rowcolor{fmblue}\color{white}Arrel & \color{white}1 & \color{white}2 & \color{white}3 & \color{white}4 & \color{white}5 & \color{white}6 & \color{white}7 & \color{white}8 & \color{white}9 & \color{white}10 & \color{white}11 & \color{white}12 & \color{white}13 & \color{white}14 & \color{white}15\\
Quadrat & 1 & 4 & 9 & 16 & 25 & 36 & 49 & 64 & 81 & 100 & 121 & 144 & 169 & 196 & 225
}\\[3pt]
{\comic\color{fmmuted}els \textbf{quadrats perfectes}: la seua arrel és exacta}
\end{center}

\vfill
\begin{minipage}[c]{0.5\linewidth}
\begin{idea}[Arrel entera i residu]
Si no és un quadrat perfecte, l'\textbf{arrel entera} és el natural més gran el quadrat del qual no passa del radicand:\\[3pt]
$\sqrt{40}$: \ $6^2=36\le40<49=7^2$ \ $\to$ \ arrel entera 6\\[2pt]
Residu: \ $40-36=4$ \qquad $40=6^2+4$
\end{idea}
\end{minipage}\hfill
\begin{minipage}[c]{0.47\linewidth}
\begin{truc}[Els zeros]
Si acaba en un nombre parell de zeros, l'arrel en té la meitat:\\[3pt]
$\sqrt{400}=20$ \qquad $\sqrt{2\,500}=50$\\[2pt]
$\sqrt{10\,000}=100$
\end{truc}
\end{minipage}

\vfill
\begin{exemple}[Exemples]
\begin{minipage}[t]{0.48\linewidth}
$\sqrt{81}=\boxed{9}$ \qquad $\sqrt{144}=\boxed{12}$ \qquad $\sqrt{900}=\boxed{30}$\\[4pt]
$\sqrt{60}$: \ $7^2=49$, $8^2=64$: \ arrel entera \boxed{7}, residu 11
\end{minipage}\hfill
\begin{minipage}[t]{0.49\linewidth}
$\sqrt{150}$: \ $12^2=144$: \ arrel entera \boxed{12}, residu 6\\[4pt]
Amb 30 rajoles, el quadrat més gran: \ $5\times5$ (25) i en sobren 5
\end{minipage}
\end{exemple}

\vfill
% ============================================================ 1.9
\clearpage
\seccio{sec9}{Problemes amb potències i arrels}{La jerarquia completa, i problemes amb quadrats i cubs.}

\vfill
\begin{minipage}[c]{0.42\linewidth}\centering
\begin{tikzpicture}[st/.style={draw=#1,fill=#1!12,line width=1.2pt,rounded corners=2mm,minimum height=0.8cm,anchor=south west,font=\large,inner xsep=6pt}]
  \node[st=fmV] at (0,2.7) {\pas{1}\ \ Parèntesis $(\ )\ [\ ]$};
  \node[st=fmA] at (0.5,1.8) {\pas{2}\ \ Potències i arrels};
  \node[st=fmK] at (1.0,0.9) {\pas{3}\ \ $\cdot$ \ i \ $:$};
  \node[st=fmB] at (1.5,0) {\pas{4}\ \ $+$ \ i \ $-$};
\end{tikzpicture}
\end{minipage}\hfill
\begin{minipage}[c]{0.55\linewidth}
\begin{exemple}[Exemple: \ el parèntesi canvia el resultat]
$2+3\cdot4^2=2+3\cdot16=2+48=\boxed{50}$\\[4pt]
$(2+3)\cdot4^2=5\cdot16=\boxed{80}$
\end{exemple}
\begin{truc}[El radical fa de parèntesi]
Primer calculem el que hi ha dins: \ $\sqrt{16+9}=\sqrt{25}=5$.
\end{truc}
\end{minipage}

\vfill
\begin{compte}[No es reparteixen en les sumes]
$\sqrt{9+16}=5$, però $\sqrt9+\sqrt{16}=7$. \qquad $(2+3)^2=25$, però $2^2+3^2=13$.
\end{compte}

\vfill
\begin{exemple}[Exemples d'operacions]
\begin{minipage}[t]{0.48\linewidth}
$20-3^2\cdot2+\sqrt{16}=20-18+4=\boxed{6}$\\[4pt]
$\sqrt{49}+2\cdot(8-3)^2=7+50=\boxed{57}$
\end{minipage}\hfill
\begin{minipage}[t]{0.49\linewidth}
$50-[18:(3^2-6)]\cdot4=50-24=\boxed{26}$\\[4pt]
$5\cdot\sqrt{36}-2^4:8=30-2=\boxed{28}$
\end{minipage}
\end{exemple}

\vfill
\begin{tcbraster}[raster columns=2,raster equal height,raster column skip=10pt,raster row skip=8pt]
\begin{exemple}[Exemple: \ un pati quadrat té 225 m$^2$. Quant fa el costat? I la tanca?]
Costat: \ $\sqrt{225}=\boxed{15\text{ m}}$\\[3pt]
Tanca: \ $4\cdot15=\boxed{60\text{ m}}$
\end{exemple}
\begin{exemple}[Exemple: \ un cub de 4 cm d'aresta, fet de cubets d'1 cm. Quants cubets té?]
$4^3=4\cdot4\cdot4=\boxed{64}$ cubets
\end{exemple}
\begin{exemple}[Exemple: \ un missatge el reben 3 persones; cada una l'envia a 3 més, i així 5 vegades]
En el 5é pas el reben \ $3^5=\boxed{243}$ persones
\end{exemple}
\begin{exemple}[Exemple: \ volem fer el quadrat més gran possible amb 90 rajoles]
$9^2=81\le90<100$: \ un quadrat de $\boxed{9\times9}$ i en sobren 9
\end{exemple}
\end{tcbraster}

\vfill
% ============================================================ resum
\resumtitol
\begin{tcbraster}[raster columns=2,raster equal height=rows,raster column skip=14pt,raster row skip=24pt,raster before skip=14pt]
\rtile{1.1}{Naturals}{$\mathbb N=\{0,1,2,\dots\}$.\\[1ex] Divisió entera: $D=d\cdot q+r$, \ $r<d$.}
\rtile{1.2}{Aproximar}{Truncar: zeros.\\[1ex] Arrodonir: mirem la xifra següent.\\[1ex] Si és 5 o més, sumem 1.}
\rtile{1.3}{Operacions}{Commutativa, associativa, distributiva.\\[1ex] Ordre: parèntesis, $\cdot$ i $:$, $+$ i $-$.}
\rtile{1.4}{Problemes}{Llig, pensa, resol i comprova.\\[1ex] Què fem amb el residu?}
\rtile{1.5}{Potències}{$a^n$: $a$ multiplicat $n$ vegades.\\[1ex] $a^1=a$, \ $a^0=1$.}
\rtile{1.6}{Base 10}{$10^n$: un 1 i $n$ zeros.\\[1ex] Descomposició polinòmica.}
\rtile{1.7}{Propietats}{$a^m\cdot a^n=a^{m+n}$, \ $a^m:a^n=a^{m-n}$,\\[1ex] $(a^m)^n=a^{mn}$.}
\rtile{1.8}{Arrels}{$\sqrt a=b\iff b^2=a$.\\[1ex] Arrel entera i residu.}
\rtile{1.9}{Jerarquia}{Parèntesis, potències i arrels, $\cdot$ i $:$, $+$ i $-$.}
\end{tcbraster}
\end{document}
